\chapter{The Code Interpreter (MicroVMs \& Firecracker)}

\begin{quote}
\textit{``A startup running a coding agent product discovered a critical breach vector in their infrastructure. Their agents were sandboxed in standard Docker containers. In a security audit, a penetration tester asked the agent to 'install the requests library and test my endpoint.' The agent ran \texttt{pip install requests} — which succeeded, because Docker containers share the host's network stack. The tester's endpoint was a data collection server. The agent then sent a POST request containing the contents of \texttt{/workspace/}, including the user's API keys stored in environment files. The Docker container had a virtual wall, but the network had an open door.''}
\end{quote}

Docker containers provide \textit{process isolation}, not \textit{machine isolation}. They share the host kernel. A Linux kernel exploit from AI-generated code does not just escape the container — it compromises the entire server, affecting every other tenant on the machine. For production agentic systems that execute arbitrary LLM-generated code, the only architecturally sound solution is hardware-level virtualization via \textbf{MicroVMs}.

\section{The Threat Hierarchy: Why Docker Is Not Enough}

\begin{center}
\begin{tabular}{|l|c|c|c|c|}
\hline
\textbf{Isolation Technology} & \textbf{Kernel Escape} & \textbf{Network Block} & \textbf{Boot Time} & \textbf{Memory Overhead} \\
\hline
No isolation & \ding{55} & \ding{55} & 0ms & 0 \\
Docker (namespaces) & Possible & Partial & 50ms & $\sim$30MB \\
gVisor (kernel intercept) & Very hard & Yes & 200ms & $\sim$100MB \\
Firecracker MicroVM & Impossible & Yes (by design) & 125ms & $\sim$5MB \\
Full QEMU VM & Impossible & Yes & 15--30s & $\sim$500MB \\
\hline
\end{tabular}
\end{center}

Firecracker is the winner: it achieves kernel-level isolation (full hardware virtualization via KVM) with near-Docker boot times and minimal memory overhead. It is deployed at scale by AWS Lambda and every major code execution service.

\section{The Firecracker Architecture in Depth}

\begin{figure}[ht]
\centering
\begin{tikzpicture}[
    scale=0.83, transform shape,
    node distance=0.6cm and 1.5cm,
    layer/.style={draw, rectangle, rounded corners, minimum width=8.5cm,
      minimum height=0.7cm, align=center, font=\small},
    sidenode/.style={draw, rectangle, rounded corners, minimum width=2cm,
      minimum height=0.65cm, align=center, font=\footnotesize},
    arrow/.style={thick,->,>=stealth}
]
\node[layer, fill=red!18] (ai_code) {AI-Generated Python Code (Untrusted)};
\node[layer, fill=orange!15, below=of ai_code] (cpython) {CPython Interpreter};
\node[layer, fill=yellow!18, below=of cpython] (guest_kernel) {Guest Linux Kernel (minimal, locked-down)};
\node[layer, fill=green!12, below=of guest_kernel] (firecracker) {Firecracker VMM (Rust process on host)};
\node[layer, fill=blue!8, below=of firecracker] (kvm) {Linux KVM (/dev/kvm, hardware-assisted)};
\node[layer, fill=gray!20, below=of kvm] (bare) {Bare Metal Hardware};

\begin{pgfonlayer}{background}
  \node[draw, dashed, thick, inner sep=0.3cm, fill=red!3,
    fit=(ai_code)(guest_kernel),
    label=right:\rotatebox{-90}{\textbf{MicroVM (Hardware Boundary)}}] {};
\end{pgfonlayer}

\node[sidenode, fill=purple!12, right=1cm of firecracker] (jailer) {Jailer\\(cgroups)};
\node[sidenode, fill=cyan!12, right=1cm of guest_kernel] (vsock) {virtio-vsock\\(host $\leftrightarrow$ guest)};
\node[sidenode, fill=green!12, left=1cm of guest_kernel] (overlayfs) {OverlayFS\\(ephemeral disk)};

\draw[arrow] (kvm) -- (firecracker);
\draw[arrow] (firecracker) -- (guest_kernel);
\draw[arrow] (guest_kernel) -- (cpython);
\draw[arrow] (cpython) -- (ai_code);
\draw[arrow, dashed] (jailer) -- (firecracker);
\draw[arrow, dashed] (vsock) -- (guest_kernel);
\draw[arrow, dashed] (overlayfs) -- (guest_kernel);
\end{tikzpicture}
\caption{Firecracker MicroVM isolation stack. The hardware boundary (KVM) means even a kernel exploit in the guest kernel cannot reach the host. All orchestrator-to-guest communication flows through the virtio-vsock, which the host fully controls.}
\end{figure}

\subsection{The Jailer: Mandatory Access Control on the VMM}
The Firecracker VMM process itself is constrained by a ``jailer'' — a companion process that applies Linux \texttt{seccomp} filters and \texttt{cgroups} \textit{before} the VMM starts. Even if Firecracker itself had a bug, the jailer's seccomp whitelist prevents the VMM from making any syscall not explicitly permitted.


\begin{center}
\noindent\rule{0.6\textwidth}{0.4pt} \\
\vspace{0.3em}
\textit{[Code implementation is available in the companion coding handbook: \texttt{coding-handbook/ch08\_code\_interpreter/}]} \\
\noindent\rule{0.6\textwidth}{0.4pt}
\end{center}


\section{OverlayFS: Sub-10ms Environment Reset}

The key to high-throughput code execution (serving thousands of concurrent agent sessions) is \textit{not} spawning a new VM for each execution. Instead, the filesystem is reset between executions using OverlayFS layering.

\begin{figure}[ht]
\centering
\begin{tikzpicture}[
    scale=0.85, transform shape,
    node distance=0.6cm and 2cm,
    layer/.style={draw, rectangle, rounded corners, minimum width=5cm,
      minimum height=0.6cm, align=center, font=\small},
    arrow/.style={thick,->,>=stealth}
]
\node[layer, fill=gray!20] (base) {Base Layer (Read-Only):\\Ubuntu + Python 3.12 + stdlib};
\node[layer, fill=blue!10, above=of base] (shared) {Shared Tools Layer (Read-Only):\\pip packages, runtime libs};
\node[layer, fill=green!15, above=of shared] (session) {Session Layer (Ephemeral Write):\\New files, mutations};
\node[layer, fill=red!10, above=of session] (code) {Agent Code Writes Here};
\node[right=1.5cm of session, align=center] (merge) {\textbf{OverlayFS merges}\\all three layers\\transparently};
\draw[arrow, dashed] (merge) -- (session);
\node[right=1.5cm of base, align=center] (reset) {\textbf{Reset}: discard session\\layer in $< 10$ms.\\Base layers untouched.};
\draw[arrow, dashed] (reset) -- (base);
\end{tikzpicture}
\caption{OverlayFS layer composition. The agent writes to the ephemeral session layer. On reset, only that layer is discarded — the base and shared layers persist across all sessions, making them instantaneously available.}
\end{figure}


\begin{center}
\noindent\rule{0.6\textwidth}{0.4pt} \\
\vspace{0.3em}
\textit{[Code implementation is available in the companion coding handbook: \texttt{coding-handbook/ch08\_code\_interpreter/}]} \\
\noindent\rule{0.6\textwidth}{0.4pt}
\end{center}


\section{Threat Vectors Inside the Sandbox}

\subsection{DNS Tunneling — Exfiltration Through Allowed UDP}
An attacker who cannot use HTTP can encode data in DNS query hostnames. DNS port 53 UDP is frequently allowed through firewalls for package installation:


\begin{center}
\noindent\rule{0.6\textwidth}{0.4pt} \\
\vspace{0.3em}
\textit{[Code implementation is available in the companion coding handbook: \texttt{coding-handbook/ch08\_code\_interpreter/}]} \\
\noindent\rule{0.6\textwidth}{0.4pt}
\end{center}


\textbf{Mitigation:} Block all outbound UDP port 53 at the hypervisor network layer. Resolve packages using an internal mirror, not the public internet.


\begin{center}
\noindent\rule{0.6\textwidth}{0.4pt} \\
\vspace{0.3em}
\textit{[Code implementation is available in the companion coding handbook: \texttt{coding-handbook/ch08\_code\_interpreter/}]} \\
\noindent\rule{0.6\textwidth}{0.4pt}
\end{center}


\subsection{CPU Exhaustion via Algorithmic Bombs}
An LLM could generate code with exponential time complexity (intentionally or accidentally). A simple nested loop can peg a CPU core indefinitely.


\begin{center}
\noindent\rule{0.6\textwidth}{0.4pt} \\
\vspace{0.3em}
\textit{[Code implementation is available in the companion coding handbook: \texttt{coding-handbook/ch08\_code\_interpreter/}]} \\
\noindent\rule{0.6\textwidth}{0.4pt}
\end{center}


\section{The VSOCK Execution Protocol}

All communication between the orchestrator and the agent's Python interpreter flows through a Virtio Socket (\texttt{AF\_VSOCK}), which is fully contained within the hypervisor — it does not touch any network interface.


\begin{center}
\noindent\rule{0.6\textwidth}{0.4pt} \\
\vspace{0.3em}
\textit{[Code implementation is available in the companion coding handbook: \texttt{coding-handbook/ch08\_code\_interpreter/}]} \\
\noindent\rule{0.6\textwidth}{0.4pt}
\end{center}

